%%%%%%%%%%%%%%%%%%%%%%%%%%%%%%%%%%%%%%%%%%%%%%%%%%%%%%%%%%%%%%%%%%%%%%%%%%%%%%%%%%
\begin{frame}[fragile]\frametitle{}
\begin{center}
{\Large Optional: For the Technically Curious}

\medskip
{\small (Feel free to mentally check out for two slides if this is not your thing, nothing after this needs it.)}
\end{center}
\end{frame}

%%%%%%%%%%%%%%%%%%%%%%%%%%%%%%%%%%%%%%%%%%%%%%%%%%%
\begin{frame}[fragile]\frametitle{The Context Window: Its Working Memory}
\begin{itemize}
\item Everything the model can ``see'' at once, your prompt, any pasted document, the conversation so far, is called the \textbf{context window}.
\item Nothing outside that window exists to the model. It does not remember yesterday's conversation unless you paste it back in.
\item This is exactly why grounding a tool in your school's own documents matters: instead of hoping the model already knows your policy, you retrieve the relevant passage and put it \emph{inside} the context window.
\item Practical rule of thumb: the more relevant material you paste in, the less the model has to guess.
\end{itemize}
\end{frame}

%%%%%%%%%%%%%%%%%%%%%%%%%%%%%%%%%%%%%%%%%%%%%%%%%%%
\begin{frame}[fragile]\frametitle{Why Hallucination Happens}
\begin{itemize}
\item A ``hallucination'' is the model producing a fluent, confident, completely fabricated answer.
\item It happens because the model is optimized to produce \emph{plausible-sounding text}, not to say ``I don't know.''
\item Ask it about a book that does not exist, and it will happily describe the plot.
\item \textbf{Two practical defenses:} give it real source material to work from (same grounding idea as before), and explicitly instruct it to say ``unsure'' when it is not confident.
\end{itemize}
\end{frame}
